\chapterdivider{12}{Scope, glossary, and technical appendix}
  {What should you remember, and where are the implementation boundaries?}
  {Reference material: terminology, derivations, engine details, and sources}

\begin{frame}{Current capability boundaries}
\small
\begin{itemize}
  \item Event detection is not transformation-safe inside \texttt{jit}/\texttt{vmap}.
  \item Event classes rely on branch-kind conventions; mixing Hopf with
        periodic data or PD/NS with equilibrium data is not rejected automatically.
  \item Dense linear/eigen algebra limits scaling to large and collocated systems.
  \item Natural continuation is intentionally unable to pass folds.
  \item Reverse-mode gradients require implicit/custom rules rather than backpropagation through the scan loop.
  \item Periodic orbits use a fixed mesh; adaptive mesh redistribution is not implemented.
  \item The periodic factory refines a supplied cycle guess but does not
        integrate to discover it; phase-plane trajectory plotting is separate.
  \item Branch switching and discovery of disconnected branches require additional analysis.
  \item Phase-plane functions are strict 2D visualization; they do not draw
        higher-dimensional slices, discover roots, classify basins, or become
        JIT/differentiable through Matplotlib.
\end{itemize}
\begin{block}{Use the tool within its scope}
JaxCont is a numerical aid for structured exploration. Combine its output with
model-specific reasoning, diagnostics, independent calculations, and domain
evidence.
\end{block}
\end{frame}

\begin{frame}{}
\begin{quote}
JaxCont follows roots $F(U,p)=0$ as a curve. For an equilibrium, $U$ is one
state. For a periodic orbit, $U$ packs the cycle samples and its period after
collocation turns the orbit into a finite root problem. Natural continuation
increments $p$ and can stall at a fold. Pseudo-arclength also corrects $p$ and
adds one geometric constraint, so its bordered system can follow the turn.

The complete predictor--corrector sweep is a bounded
\texttt{lax.while\_loop} over fixed buffers and is JIT-compiled. JAX supplies
Jacobians, batching and accelerator execution; a custom VJP supplies implicit
root sensitivities. Equilibrium stability uses Jacobian eigenvalues, while
cycle stability uses Floquet multipliers. These tools accelerate and organize
the analysis, but a study still needs a valid start, diagnostics and
independent validation. For a 2D autonomous model, the same problem can also
be frozen at one parameter to connect the branch diagram with nullclines,
continued equilibria, vector fields and trajectories.
\end{quote}
\end{frame}

\begin{frame}{Vocabulary checklist}
\scriptsize
\begin{description}
  \item[Continuation] following a connected family of solutions.
  \item[Predictor] a local guess for the next branch point.
  \item[Corrector] Newton iterations that return the guess to the branch.
  \item[Fold] a regular turning point where $F_{\state}$ is singular.
  \item[Pseudo-arclength] continuation using an approximate curve-distance coordinate.
  \item[Bordered system] the Jacobian enlarged with parameter and constraint rows/columns.
  \item[Collocation] enforcing an ODE at selected points of a piecewise-polynomial orbit.
  \item[Phase condition] one equation that removes arbitrary time shift of a cycle.
  \item[Floquet multiplier] eigenvalue of the one-period linearized map.
  \item[Monodromy matrix] the linear map advancing perturbations by one period.
  \item[Nullcline] a curve where one component of a 2D vector field is zero.
  \item[Phase plane] the true state-space flow of a two-state autonomous
        system at one fixed parameter.
  \item[JIT] trace and compile array code for later reuse.
  \item[\texttt{lax.while\_loop}] JAX's compiled, traceable dynamic loop primitive.
  \item[\texttt{vmap}] automatic batching of one function across an array axis.
  \item[PyTree] a nested data structure JAX knows how to flatten into array leaves.
\end{description}
\end{frame}

\begin{frame}{Main takeaways}
\begin{enumerate}
  \item One root-based abstraction now covers both equilibria and collocated periodic orbits.
  \item Pseudo-arclength and the bordered Newton system follow ordinary folds.
  \item Stability is problem-specific: $\operatorname{Re}\lambda<0$ for
        equilibria, $|\mu|<1$ for nontrivial Floquet multipliers.
  \item JAX makes derivatives, whole-loop compilation, batching and implicit
        sensitivities composable in one numerical codebase.
  \item Fixed buffers, static mesh shapes and eager-only events are deliberate
        transformation trade-offs.
  \item For 2D models, phase planes connect branch and stability results to
        nullclines and trajectories without changing the solver core.
  \item A trustworthy study begins with a verified solution and ends with
        residual, resolution and independent validation.
\end{enumerate}
\end{frame}

\begin{frame}{Appendix: natural versus pseudo-arclength}
\small
\begin{center}
\begin{tabular}{p{.25\textwidth}p{.32\textwidth}p{.32\textwidth}}
\toprule
 & Natural & Pseudo-arclength\\
\midrule
Unknown during correction & $\state$ & $(\state,p)$\\
Fixed quantity & predicted $p$ & approximate distance from previous point\\
Newton matrix & $F_{\state}$ & $\begin{bmatrix}F_{\state}&F_p\\\dot\state^T&\dot p\end{bmatrix}$\\[3mm]
At ordinary fold & singular / stalls & usually remains regular\\
Cost per Newton step & $n\times n$ solve & $(n+1)\times(n+1)$ solve\\
JaxCont engine & \texttt{natural\_scan} & \texttt{pseudo\_arclength\_scan}\\
\bottomrule
\end{tabular}
\end{center}
\end{frame}

\begin{frame}{Appendix: scalar cubic fold derivation}
For
\[
F(u,p)=p+u-u^3/3=0,
\]
the branch can be written
\[
p(u)=u^3/3-u.
\]
Folds occur when
\[
\frac{dp}{du}=u^2-1=0
\quad\Longrightarrow\quad u=\pm1.
\]
Therefore
\[
(u,p)=(-1,2/3),\qquad(1,-2/3).
\]
Also $F_u=1-u^2$, so the same points have a zero scalar Jacobian. This example connects the geometry, the singularity and the detector's expected value.
\end{frame}

\begin{frame}[fragile]{Appendix: simplified engine pseudocode}
\begin{lstlisting}
allocate fixed buffers
tangent = tangent_at(initial_point)

while attempts < max_steps and not stop:
    predicted = accepted + ds * tangent
    candidate, ok, nit = bordered_newton(predicted)
    candidate_tangent = tangent_at(candidate)

    buffers = conditional_write(buffers, candidate, ok)
    accepted = where(ok, candidate, accepted)
    tangent = where(ok, candidate_tangent, tangent)
    ds = adapt(ds, nit, ok)
    stop = reached_target or stalled or nonfinite

return fixed_buffers, n_valid
\end{lstlisting}
The real code uses JAX arrays, \texttt{NamedTuple} carry state, \texttt{jnp.where}, and nested \texttt{lax.while\_loop}s.
\end{frame}

\begin{frame}{Appendix: sources used for this deck}
\small
\begin{itemize}
  \item implementation in \texttt{src/jaxcont/core/scan\_continuation.py}
  \item public dispatch and result paths in \texttt{src/jaxcont/api.py}
  \item collocation and periodic problem construction in
        \texttt{core/collocation.py} and \texttt{problems/periodic.py}
  \item Floquet and event logic in \texttt{stability/floquet.py} and
        \texttt{bifurcations/events.py}
  \item phase-plane composition in \texttt{viz/phase\_plane.py}, built on
        \texttt{BifProblem.as\_rhs(p)}
  \item v0.2 and phase-plane design specs and implementation plans in
        \texttt{docs/superpowers/}
  \item tests and Examples 04, 06--12 for folds, batching, differentiation,
        periodic continuation, PD/NS and FitzHugh--Nagumo phase planes
\end{itemize}
\begin{block}{Reproducibility note}
Re-run numerical examples and benchmarks with the current implementation,
settings, precision, and target hardware before external publication.
\end{block}
\end{frame}
